\steppingstone{What Happens When We Join Two Collections?}
\label{stone:adding}

In the previous stepping stone we learned that a PhiBase quantity records an embedding. The pair

\[
P(n,p)
\]

does not describe the order of the pieces. It records only the total number of short units and the total number of long units contained in the construction.

That observation makes addition almost unavoidable.

Suppose one collection contains

\[
P(3,1),
\]

and a second contains

\[
P(2,4).
\]

Imagine placing all of the pieces on the same workbench.

Nothing has changed.

The short pieces are still short.

The long pieces are still long.

The only question is how many of each kind are now present.

There are

\[
3+2=5
\]

short units and

\[
1+4=5
\]

long units.

The combined collection is therefore

\[
P(5,5).
\]

Notice what did \emph{not} matter.

We never asked which piece came first.

We never asked which piece came last.

Order belongs to a path.

Addition belongs to a collection.

The embedding simply counts the total amount of each kind of length.

The same construction works for every pair of PhiBase quantities.

If one collection is

\[
P(a,b)
\]

and the other is

\[
P(c,d),
\]

their combined embedding is

\[
P(a,b)+P(c,d)
=
P(a+c,\;b+d).
\]

Nothing mysterious has happened.

The short counts combine with the short counts.

The long counts combine with the long counts.

That is exactly what the construction itself tells us to do.

Subtraction follows the same reasoning.

If a collection described by

\[
P(8,7)
\]

contains another collection described by

\[
P(1,5),
\]

removing those pieces leaves

\[
P(8,7)-P(1,5)=P(7,2).
\]

Once again, only the totals change.

The embedding remembers how much material remains.

It says nothing about the order in which the pieces were originally assembled.

That distinction becomes important later.

When we begin walking through a construction, the order of the steps creates an address.

Today we are not walking.

We are only counting.

\section*{Builder's Notebook}

Combine the following collections.

\[
P(2,3)+P(5,1),
\]

\[
P(8,0)+P(0,4),
\]

\[
P(6,5)-P(2,1),
\]

\[
P(9,7)-P(4,6).
\]

Now write two different arrangements of short and long pieces that both represent

\[
P(4,3).
\]

The arrangements are different.

The embedding is the same.

\section*{Looking Ahead}

Adding collections changes only the totals.

The next question is very different.

What happens when every short piece grows into a long one, and every long piece grows by the same amount?
